\section{Background and Related Work}\label{sec:background}

\subsection{Blockchain Foundations for Health Surveillance}\label{subsec:blockchain}

Blockchain provides four properties directly relevant to distributed health surveillance: immutability of committed records, decentralized validation without trusted intermediaries, consensus-driven Byzantine fault tolerance, and end-to-end transaction traceability.\cite{yaga2018blockchain} However, these properties alone are insufficient for privacy-sensitive, resource-constrained surveillance deployments. Full ledger replication forces every node to store the complete surveillance payload, producing unsustainable storage growth and bandwidth demands in low- and middle-income settings. Immutability, while protecting against retroactive tampering, does not prevent metadata exposure, overbroad read access across health-system tiers, or the absence of selective disclosure. Consequently, modern blockchain-health architectures have shifted from monolithic on-chain designs toward \emph{hybrid} models that confine the ledger to coordination artifacts---hashes, permissions, provenance, alerts, and audit events---while encrypted bulk data resides off-chain.

\subsection{Platform and Consensus Rationale}\label{subsec:platform}

Ethereum was selected as the coordination layer based on its native smart-contract support, Turing-complete programmability (Solidity), and mature tooling for permissioned deployment (Table~\ref{tab:platform}). Critically, the architecture reframes Ethereum from a bulk data repository to a \emph{selective-anchoring} layer: smart contracts manage only cryptographic hashes, access-control state, and audit artifacts, while encrypted surveillance payloads are stored off-chain via IPFS. This separation addresses the scalability, cost, and privacy limitations inherent in direct Layer-1 data recording.

\begin{table}[!ht]
\caption{Blockchain platform comparison (selection criteria).}\label{tab:platform}
\begin{tabular*}{\columnwidth}{@{\extracolsep\fill}lccc@{}}
\toprule
Criterion & Ethereum & Bitcoin & XRPL \\
\midrule
Smart contracts & Yes & No & No \\
dApp support & Yes & No & No \\
Turing-complete language & Yes & No & No \\
Digital identity & Yes & Yes & No \\
Data oracles & Yes & Yes & No \\
Human-readable addresses & Yes & Yes & No \\
Block time & 14 s & 10 min & 4 s \\
\bottomrule
\end{tabular*}
\end{table}

Consensus selection is treated as a workload-specific governance decision rather than an assumed default. While the prototype employs Proof of Authority (PoA) for its low latency and deterministic finality, Section~\ref{sec:model} presents a comparative evaluation of PoA against Delegated Proof of Stake (DPoS), Practical Byzantine Fault Tolerance (PBFT) variants, and Raft-based alliance models, assessed against LMIC workload characteristics, trust topology, validator appointment, collusion resilience, and infrastructure constraints.

\subsection{Disease Early Warning and Surveillance in Ethiopia}\label{subsec:dews}

Ethiopia's DEWS system operates through a five-tier hierarchy: community health posts, woreda health offices, zone health offices, regional health offices, and the national Ethiopian Public Health Institute (EPHI)/Public Health Emergency Management (PHEM) center. The system conducts indicator-based surveillance (structured scheduled reporting) and event-based surveillance (rapid informal signal capture). Twenty priority diseases are classified into 13 immediately reportable conditions (notification within 30 minutes) and 7 weekly reportable conditions, with report completeness calculated weekly at every tier. Laboratory surveillance operates through a parallel tiered network (woreda, zone, regional, national) responsible for specimen confirmation.

This administrative hierarchy maps directly to the proposed multi-tier computing architecture. Peripheral health posts and community workers operate in environments characterized by intermittent connectivity, underpowered devices, and limited technical capacity---conditions that preclude full blockchain peer responsibilities. District and zonal offices possess intermediate infrastructure suitable for batching, validation, and edge/fog preprocessing. Regional and national tiers command sufficient resources for validator duties and analytical workloads. The architecture therefore assigns asymmetric node roles aligned to each tier's actual capacity, ensuring that resource-constrained peripheral sites participate as lightweight reporting nodes rather than heavy consensus participants.

\subsection{Related Work}\label{subsec:related}

\subsubsection{Blockchain-Based Healthcare Architectures}

Blockchain adoption in healthcare has progressed through identifiable architectural generations. Early systems stored clinical records directly on-chain: Omar et al.\cite{omar2019privacy} proposed a patient-centric blockchain using decentralization to mitigate cloud security risks, but lacked interoperability between actors and placed full payloads on the ledger. Griggs et al.\cite{griggs2018healthcare} applied smart contracts to medical IoT data with HIPAA-compliant notifications, yet reported transaction delays that limit real-time utility. Quaini et al.\cite{quaini2018model} presented UniRec, a distributed electronic health record architecture whose full-replication design produces prohibitive disk space requirements at scale. These first-generation approaches share a common limitation: treating the blockchain as a bulk data repository rather than a coordination layer, thereby inheriting the storage, cost, and privacy weaknesses identified in Section~\ref{sec:introduction}.

More recent architectures have moved toward hybrid storage and privacy-layered designs. Zhang et al.\cite{zhang2018fhirchain} proposed FHIRChain, combining blockchain coordination with HL7 FHIR interoperability for collaborative cancer care, but reported semantic interoperability challenges and legacy system limitations. The architectural direction established by such hybrid approaches---separating data custody from on-chain coordination---informs the selective-anchoring model proposed in this work, which extends it with encrypted off-chain storage via IPFS, zero-knowledge verification, and dynamic permission lifecycle management.

\subsubsection{Surveillance-Specific Digital Architectures}

Digital surveillance systems have predominantly employed machine learning. Bollig et al.\cite{bollig2020machine} applied LSTM networks to syndromic surveillance across 33,000 veterinary necropsy reports, demonstrating temporal and spatial pattern detection. Harvey et al.\cite{harvey2021predicting} developed Gaussian Process and Random Forest models for malaria early warning in Burkina Faso using routine health facility data. Rubio-Solis et al.\cite{rubiosolis2019zika} proposed neural network-based predictive surveillance for dengue at district level in Brazil. While these systems demonstrate predictive utility, they require centralized data aggregation, remain vulnerable to adversarial manipulation, and provide no inherent guarantees of provenance, immutability, or audit traceability---properties essential for trustworthy multi-institutional surveillance.

\subsubsection{Architectural Gap Analysis}

Table~\ref{tab:related} positions the present work against prior systems across the six architectural dimensions identified in the blueprint.

\begin{table*}[!ht]
\caption{Comparative analysis of related work against target architectural dimensions.}\label{tab:related}
\begin{tabular*}{\textwidth}{@{\extracolsep\fill}llllll@{}}
\toprule
Study & Hybrid Storage & Layered Privacy & Edge/Fog & Consensus Rationale & LMIC Awareness \\
\midrule
Omar et al.\cite{omar2019privacy} & No & No & No & No & No \\
Griggs et al.\cite{griggs2018healthcare} & No & Partial & No & No & No \\
Quaini et al.\cite{quaini2018model} & No & No & No & No & No \\
Zhang et al.\cite{zhang2018fhirchain} & Partial & Partial & No & No & No \\
Bollig et al.\cite{bollig2020machine} & N/A & No & No & N/A & No \\
Harvey et al.\cite{harvey2021predicting} & N/A & No & No & N/A & Partial \\
Rubio-Solis et al.\cite{rubiosolis2019zika} & N/A & No & No & N/A & No \\
\textbf{This work} & \textbf{Yes} & \textbf{Yes} & \textbf{Yes} & \textbf{Yes} & \textbf{Yes} \\
\bottomrule
\end{tabular*}
\end{table*}

\textbf{Research gap.} No prior blockchain-health architecture integrates all five dimensions required for LMIC disease surveillance: (1)~hybrid storage with selective on-chain anchoring, (2)~layered privacy with zero-knowledge verification and dynamic permission lifecycle, (3)~edge/fog preprocessing for resource-constrained peripheral sites, (4)~comparative consensus evaluation with explicit validator governance, and (5)~LMIC-aware asymmetric node-role design aligned to public-health administrative hierarchies. The present work addresses this compound gap by proposing a unified architecture that treats privacy, consensus, scalability, and resource asymmetry as linked governance concerns within a single hierarchical surveillance framework.

%% ===== RESEARCH METHODOLOGY =====
